\exercisesfor{8}

设 \(p\) 是域 \(K\) 的剩余特征，且 \(>0\)。记 \(\mathfrak{o}_{K}\) 为赋值 \(v\) 在域 \(K\) 上的环。

\begin{enumerate}
\def\labelenumi{\arabic{enumi}.}
\tightlist
\item
  \begin{enumerate}
  \def\labelenumii{(\alph{enumii})}
  \tightlist
  \item
    若元素 \(\pi\) 属于 \(K\) 且满足 \(v(\pi) = \theta\) 以及 \(\pi^{p-1} / p \equiv -1 \pmod{\pi o_K}\)，则称其为可容许元素。
  \end{enumerate}
\end{enumerate}

（注意 \(\pi^{p-1}/p\in\mathfrak{o}_{K}\)，因为 \(v(\pi)=\theta.\)）

证明存在满足本段条件且含有可容许元素的域 K（向 \(Q_{p}\) 添加一个 \((p - 1)\) 次根 \(-p\)）。

\begin{enumerate}
\def\labelenumi{(\alph{enumi})}
\setcounter{enumi}{1}
\tightlist
\item
  令 \(\pi\) 是 K 中的可容许元素。考虑下列系数属于 K 的形式幂级数：
\end{enumerate}

\[
e _ {\pi} (\mathrm{U}) = \frac {1}{\pi} e (\pi \mathrm{U}) = \sum_ {n = 1} ^ {\infty} \pi^ {n - 1} \mathrm{U} ^ {n} / n!,
\]

\[
l _ {\pi} (\mathrm{U}) = \frac {1}{\pi} l (\pi \mathrm{U}) = \sum_ {n = 1} ^ {\infty} (- \pi) ^ {n - 1} \mathrm{U} ^ {n} / n.
\]

证明这些级数互为逆级数，且其系数属于 \(\mathfrak{o}_{K}\)（使用引理 2）。

\begin{enumerate}
\def\labelenumi{(\alph{enumi})}
\setcounter{enumi}{2}
\tightlist
\item
  令 \(\{\mathrm{U},\mathrm{V}\}\) 是含两个元素的集合，令 \(\mathrm{H}(\mathrm{U},\mathrm{V})\) 为 Hausdorff 级数。记
\end{enumerate}

\[
\mathrm{H} _ {\pi} (\mathrm{U}, \mathrm{V}) = \frac {1}{\pi} \mathrm{H} (\pi \mathrm{U}, \pi \mathrm{V}).
\]

于是 \(H_{\pi}(U, V) \in \hat{L}_{K}(\{U, V\})\)。

证明

\[
\mathrm{H} _ {\pi} (\mathrm{U}, \mathrm{V}) = l _ {\pi} (e _ {\pi} (\mathrm{U}) + e _ {\pi} (\mathrm{V}) + \pi e _ {\pi} (\mathrm{U}). e _ {\pi} (\mathrm{V})).
\]

推出 \(H_{\pi} \in \hat{L}_{\mathfrak{o}_{K}}(\{\mathrm{U}, \mathrm{V}\})\)，即 \(H_{\pi}\) 的系数属于 \(o_{K}\)；从而得到命题 1 的另一种证明。

\begin{enumerate}
\def\labelenumi{(\alph{enumi})}
\setcounter{enumi}{3}
\tightlist
\item
  令 \(\tilde{e}_{\pi}, \tilde{l}_{\pi}\) 和 \(\tilde{\mathbf{H}}_{\pi}\) 表示将 \(e_{\pi}, l_{\pi}\) 和 \(\mathbf{H}_{\pi}\) 模 \(\pi \mathfrak{o}_{K}\) 化简所得的级数；它们的系数属于环 \(\mathfrak{o}_{K} / \pi \mathfrak{o}_{K}\)。
\end{enumerate}

证明 \(\tilde{l}_{\pi}(\mathrm{U}) = \mathrm{U} - \mathrm{U}^{p}\)（注意 \(v_{p}(n) < (n - 1)\theta\) 对 \(n \neq 1, p\) 成立）。推出

\[
\tilde {e} _ {\pi} (\mathrm{U}) = \mathrm{U} + \mathrm{U} ^ {p} + \dots + \mathrm{U} ^ {p ^ {n}} + \dots
\]

并且

\[
\tilde {\mathrm{H}} _ {\pi} (\mathrm{U}, \mathrm{V}) = \tilde {e} _ {\pi} (\mathrm{U}) + \tilde {e} _ {\pi} (\mathrm{V}) - (\tilde {e} _ {\pi} (\mathrm{U}) + \tilde {e} _ {\pi} (\mathrm{V})) ^ {p},
\]

或者也可写为

\[
\tilde {\mathrm{H}} _ {\pi} (\mathrm{U}, \mathrm{V}) = \mathrm{U} + \mathrm{V} - \Lambda_ {p} \left(\sum_ {n = 0} ^ {\infty} \mathrm{U} ^ {p ^ {n}}, \sum_ {n = 0} ^ {\infty} \mathrm{V} ^ {p ^ {n}}\right),
\]

其中 \(\Lambda_{p}\) 由下式定义：

\[
\Lambda_ {p} (\mathrm{U}, \mathrm{V}) = (\mathrm{U} + \mathrm{V}) ^ {p} - \mathrm{U} ^ {p} - \mathrm{V} ^ {p}, \quad \text { cf.   Chapter   I,   } \S 1, \text {   Exercise   19. }
\]

特别地，\(\tilde{\mathrm{H}}_{\pi}(\mathrm{U},\mathrm{V})\) 属于 \(\hat{\mathrm{L}}_{\mathbf{F}_{p}}(\{\mathrm{U},\mathrm{V}\})\)，且其 \(p\) 次齐次分量为 \(-\Lambda_p(\mathrm{U},\mathrm{V})\)。于是

\[
\begin{array}{c} \tilde {\mathrm{H}} _ {\pi} (\mathrm{U}, \mathrm{V}) = \tilde {\mathrm{H}} _ {\pi} (\mathrm{V}, \mathrm{U}) \\ \tilde {\mathrm{H}} _ {\pi} (\mathrm{U}, \tilde {\mathrm{H}} _ {\pi} (\mathrm{V}, \mathrm{W})) = \tilde {\mathrm{H}} _ {\pi} (\tilde {\mathrm{H}} _ {\pi} (\mathrm{U}, \mathrm{V}), \mathrm{W}) \quad \text { in } \hat {\mathrm{L}} _ {\mathbf {F} _ {p}} (\{\mathrm{U}, \mathrm{V}, \mathrm{W} \}), \end{array}
\]

其中 W 是第三个不定元。

\begin{enumerate}
\def\labelenumi{(\alph{enumi})}
\setcounter{enumi}{4}
\tightlist
\item
  令 C(U, V) = H(-U, H(-V, H(U, V)))（“Hausdorff 交换子”）。于是 exp(C(U, V)) = exp(-U) exp(-V) exp(U) exp(V)。记
\end{enumerate}

\[
\mathrm{C} _ {\pi} (\mathrm{U}, \mathrm{V}) = \frac {1}{\pi} \mathrm{C} (\pi \mathrm{U}, \pi \mathrm{V}).
\]

证明 \(C_{\pi}(U,V)\) 的系数属于 \(\pi o_{K}\)（使用 \(\tilde{H}_{\pi}(U,V)=\tilde{H}_{\pi}(V,U)\) 这一事实）。†

\begin{enumerate}
\def\labelenumi{\arabic{enumi}.}
\setcounter{enumi}{1}
\tightlist
\item
  我们知道（§ 6，习题 3），若 \(n \geqslant 2\)，则 H(U, V) 的双次数为 \((n, 1)\) 的分量是 \(\frac{1}{n!} b_n(\mathrm{ad}\mathrm{U})^n (\mathrm{V})\)，其中 \(b_n\) 是第 \(n\) 个伯努利数。
\end{enumerate}

由此和命题 1 推出不等式

\[
v _ {p} (b _ {n} / n!) \leqslant n / (p - 1).
\]

利用克劳森—冯·施陶特定理（《实变函数》，第六章，§2，习题 6）重新得到这一结果，并证明等号成立当且仅当 \(S(n) = p - 1\)。

\begin{enumerate}
\def\labelenumi{\arabic{enumi}.}
\setcounter{enumi}{2}
\tightlist
\item
  假设 K 含有本原 \(p\) 次单位根 \(w\)。记 \(\pi = w - 1\)。利用公式
\end{enumerate}

\[
w ^ {i} - 1 = \pi (1 + w + \dots + w ^ {i - 1})
\]

证明 \(v(w^{i} - 1) = v(\pi)\) 对 \(1 \leqslant i \leqslant p - 1\) 成立，并且

\[
\frac {w ^ {i} - 1}{\pi} \equiv i \pmod {\pi}.
\]

利用公式 \(p = \prod_{i=1}^{p-1} (w^i - 1)\) 推出 \(\pi\) 是可容许元素（习题 1）。

¶ 4. 令 c 是满足 \(\geqslant1\) 的整数，令 P 是包含所有 \(\leqslant c\) 的素数的素数集合。令 \(Z_{P}=S_{P}^{-1}Z\)（§ 4，习题 16）。证明 Hausdorff 级数 H(U, V) 中次数 \(\leqslant c\) 的项属于 \(L_{Z_{P}}(\{U, V\})\)。推出，若 g 是幂零李 \(Z_{P}\)-代数，且类为 \(\leqslant c\)，则复合律 \((u,v)\mapsto\mathrm{H}(u,v)\) 使 g 成为类 \(\leqslant c\) 的 P-无挠、P-可除幂零群。证明反过来每个具有这些性质的群都可由此过程得到。（使用“Hausdorff 公式的逆”，参见~§6，习题 4，并证明其中出现的指数 \(\alpha(m)\)、\(\beta(m)\) 属于 \(Z_{P}\)，当 \(l(m)\leqslant c\) 时。）

特别地，每个阶为 \(p^{n}\) 且类 \textless p 的 p-群，都可以通过 Hausdorff 律由一个类 \textless p 且含有 \(p^{n}\) 个元素的幂零李 Z-代数得到。（取 P 为异于 p 的素数集合。）

